\begin{abstract}
A large literature scores citation correctness as a single signal. We
decompose it into three operationally independent constraints:
whether the cited span \textbf{exists} in the corpus, whether it was
\textbf{in-context} (in the retrieval set the model was shown for this
query), and whether it \textbf{supports} the claim. We externalize
each as a validator gate that refuses output rather than degrading it,
and build a benchmark of 254 FDA-drug-label seeds expanded to 1{,}270
examples via named matched-pair failure-injection operators. On the v2
injected FDA-label benchmark --- a controlled benchmark, not a
prevalence study --- the structural rules (existence, in-context) are
exact and model-independent (F1 = 1.000); the three rules catch fully
disjoint classes of error (100\% by the operational blindness
measure); and the composed validator improves should-refuse F1 by
$+56$ points over the best non-composed baseline. The semantic support
gate, measured separately, catches genuine failures at $0.99$ recall,
but a strong NLI model still conservatively over-refuses true
claims --- a finding we measure rather than hide, and which motivates
treating support as its own refuse-or-resolve gate. We release the
system, the benchmark, and one-command reproduction.

\paragraph{Scope.} All results are on a controlled, failure-injected,
single-domain (FDA drug-label) benchmark; they characterize validator
behavior under matched-pair injected failures, not real-world failure
prevalence. Structural rules (existence, in-context) are exact and
model-independent; the support gate is a conservative, high-recall
safety gate, not a high-precision throughput filter.
\end{abstract}

\section{Introduction}
\label{sec:intro}

Consider a model that writes \emph{``the trial enrolled 1000 patients''}
and attaches a citation to a span reading \emph{``the trial enrolled 240
patients.''} The citation is well-formed, the link resolves, and the
cited span is topically relevant --- yet the claim is false. Most
existing attribution metrics score this as ``cited with high
relevance.'' That conflation is the problem this paper addresses.

``Is this citation correct?'' is treated as one signal, but it is
three. A citation is trustworthy only if (1) the cited span
\textbf{exists}, (2) it was \textbf{in the context} the model was
given for this query, and (3) it \textbf{supports} the claim. These
fail independently: a fabricated span ID fails (1); a real span the
model was never shown fails (2) while passing (1); a real, shown, but
irrelevant span passes (1) and (2) and fails (3). Conflating them into
one score hides the practitioner's actual question --- \emph{which
rule failed?} --- and the three failures have completely different
fixes (decoding/format, retrieval-context, reasoning).

\paragraph{Contributions.}
(i) We decompose citation correctness into existence, in-context, and
support, and externalize each as a validator gate that refuses rather
than degrades. (ii) We build a public FDA-drug-label benchmark with
controlled, named matched-pair failure injection (254 seeds $\to$
1{,}270 labeled examples). (iii) We show, under matched-pair failure
injection, that the three rules catch \textbf{non-overlapping} classes
of error (fully disjoint on the v2 injected FDA-label benchmark, 100\%
by the operational blindness measure) and that the composed validator
improves should-refuse F1 by \textbf{$+56$ points} over the best
non-composed baseline; the structural rules are exact and
model-independent.

Section~\ref{sec:related} places the work; Section~\ref{sec:rules}
defines the rules; Section~\ref{sec:bench} the benchmark;
Section~\ref{sec:exp} the experiments; Section~\ref{sec:limits}
limitations; Section~\ref{sec:conc} conclusion.

\section{Related Work}
\label{sec:related}

\paragraph{Attribution and faithfulness.}
The dominant operational definition of attribution is AIS
\citep{rashkin2023ais}, which judges whether a model statement is
fully derivable from an identified source via a two-stage human
protocol (interpretability, then an attribution rating). AIS bundles
all three of our constraints into a single human judgment: the source
is shown to the rater, so existence and in-context grounding are
satisfied tautologically and only support is genuinely assessed.
Attributed QA \citep{bohnet2022attributedqa} defines the
attributed-answering task and benchmarks retrieve-then-generate
variants, but reuses AIS for evaluation and inherits its bundling.
TRUE \citep{honovich2022true} standardizes factual-consistency
evaluation and finds large-scale NLI and question-generation-and-%
answering to be strong and complementary; it establishes the
NLI-for-faithfulness lineage our support gate draws on but is
orthogonal to attribution decomposition.

\paragraph{RAG citation benchmarks.}
ALCE \citep{gao2023alce} is the standard generation-with-citations
benchmark; its redundancy-aware citation precision measures support
carefully via an NLI model, but citations are integer indices into the
retrieved set by construction ($c_{i,j} \in D$), so existence and
in-context grounding are \emph{assumed true by construction and cannot
be measured}. HAGRID \citep{kamalloo2023hagrid} reuses the ALCE-style
metric on a human-LLM collaborative corpus and carries the same
single-signal limitation. ExpertQA \citep{malaviya2024expertqa}
provides expert-curated questions with multi-attribute evaluation, but
attribution is one bundled dimension among several and is not itself
decomposed into existence, in-context, and support.

\paragraph{Span-level and constrained citing.}
GopherCite \citep{menick2022gophercite} trains a model to answer with
verbatim quotes and \emph{enforces} quote-document fidelity through
constrained decoding rather than measuring it: existence and
in-context grounding cannot fail by construction, and its primary
``Supported \& Plausible'' metric bundles support with plausibility.
It shows the field has groped toward gate-1/gate-2 enforcement without
naming it as a separable, measurable concern.

\paragraph{Verifiability and post-hoc revision.}
\citet{liu2023verifiability} audit commercial generative search
engines and report that only about half of generated sentences are
fully supported by their citations; it is closest in spirit to this
work but notes URL non-resolution only in passing and cannot measure
in-context grounding, since deployed search engines do not expose a
retrieved set. RARR \citep{gao2023rarr} post-edits model output to fix
unsupported content, operating on support with revise-the-output
semantics --- the inverse of our refuse-or-resolve deployment
philosophy. Self-RAG \citep{asai2024selfrag} is the closest existing
decomposition: it trains reflection tokens that separately rate
passage relevance (\texttt{IsRel}) and three-valued support
(\texttt{IsSup}). But \texttt{IsRel} measures retrieval
\emph{relevance}, not whether a cited span was in the prompt for
\emph{this} query, so it addresses neither existence nor our
in-context gate; and both signals are \emph{internalized} trained
tokens, not \emph{externalized} validator gates that refuse the
output.

\paragraph{Frameworks and methodology.}
RAGAS \citep{es2024ragas} is a reference-free RAG evaluation framework
whose faithfulness metric is essentially our support gate via
LLM-as-judge, while its context-relevance metric evaluates retrieval
rather than citation grounding; we position the three-rule
decomposition as the right granularity for RAGAS-shaped tools to
adopt. The synthetic failure-injection design follows the contrast-set
methodology of \citet{gardner2020contrast}. Recent agentic-citation
work is adjacent but disjoint: CiteGuard \citep{choi2025citeguard}
reframes citation evaluation as predicting which academic paper a
human would cite (the CiteME task), a recommendation problem rather
than grounding validation; and accounts of citation hallucination as
post-rationalization rather than genuine reliance
\citep{wallat2024correctness} operate on the model's reasoning to
explain \emph{why} failures occur, complementary to our
validator-boundary question of \emph{which} failure occurred.

\paragraph{The gap.}
No prior work separates existence, in-context, and support as three
independently \emph{measured} gates at the validator boundary, and
none demonstrates that they catch distinct, non-overlapping classes of
citation errors. The closest near-misses either bundle the three into
one judgment (AIS), assume two of them away by construction (ALCE,
GopherCite), internalize a partial decomposition as trained signals
(Self-RAG), or study a setting where in-context grounding is
unmeasurable \citep{liu2023verifiability}. This paper measures all
three as external, refusing gates and shows their catch sets are
disjoint on the benchmark. We claim citation \emph{correctness} at the
validator boundary, not causal faithfulness: the in-context gate shows
a span was \emph{available} to the model, not that the model
\emph{relied} on it, and post-rationalization remains possible
\citep[cf.][]{wallat2024correctness} --- causal reliance is a
separate, out-of-scope dimension.


\section{The Three Rules}
\label{sec:rules}

\subsection{Formal definitions}
\label{sec:rules-formal}

Let $C$ be the corpus of spans. For a query $q$, let $R_q \subseteq C$
be the retrieval set the model was shown. Let $s$ be a cited span and
$y$ the claim sentence. The three gates:
\begin{align*}
  \mathrm{exists}(s, C) &\equiv s \in C, \\
  \mathrm{in\_context}(s, R_q) &\equiv s \in R_q, \\
  \mathrm{supports}(s, y) &\equiv \mathrm{NLI}(\mathrm{text}(s), y) = \mathrm{entailment}.
\end{align*}
The validator accepts the citation $(s, y)$ iff
$\mathrm{exists} \land \mathrm{in\_context} \land \mathrm{supports}$,
and otherwise refuses with the typed reason of the first failing gate.
$\mathrm{in\_context}$ establishes that $s$ was \emph{available} to
the model, not that the model \emph{used} it (see Scope of the claim,
below).

\subsection{Existence}
The cited identifier must resolve to a real span in the corpus.
Failure: a fabricated pointer into nothing. Implementation: a
primary-key lookup, $O(1)$, deterministic.

\subsection{In-context}
The cited span must have been in the retrieval set assembled for
\emph{this} query --- the closed-loop constraint: a model may only
cite what it was shown. Failure: citing a real span the model was not
given this turn (back-filled from parametric memory). Implementation:
a per-prompt allowed-span \texttt{HashSet}, deterministic.

\subsection{Support}
The cited span must entail the claim. Implementation: an NLI
cross-encoder over \texttt{(span\_text, sentence)} pairs producing a
three-class verdict (entailment / neutral / contradiction). The only
model-dependent gate.

\subsection{Composition and strict-wins}
The gates run in order --- existence, then in-context, then
support --- so each cleanly attributes one failure class (a class is
owned by the first gate that can see it). Support aggregation is
\textbf{strict-wins}: any contradiction wins, else any neutral wins,
else support. ``Every cited span must hold on its own'' is a
conservative safety stance; a concatenated-evidence alternative is
noted as future work, not shipped.

\subsection{Validator output and deployment}
Refusal types: \texttt{Uncited}, \texttt{UnknownSpan} (existence),
\texttt{OutOfContext} (in-context), \texttt{Unsupported},
\texttt{Contradicted} (support). Deployment philosophy is
\textbf{refuse-or-resolve} --- never silently degrade. A conservative
gate is the correct default when the cost of a confidently-wrong
citation is high (the target domain is pharmaceutical text).

\subsection*{Scope of the claim: citation correctness, not causal faithfulness}

We validate citation \emph{correctness} at the output boundary. The
in-context gate verifies the cited span was in the retrieval set the
model was shown; it does \textbf{not} verify the model causally relied
on that span --- a model can post-rationalize a well-formed citation.
Causal reliance (counterfactual span removal, attribution tracing) is
a distinct fourth dimension, explicitly out of scope. This boundary is
deliberate: existence, in-context, and support are checkable at the
validator without model internals; causal faithfulness is not. Recent
work separates citation correctness from faithfulness and reports
substantial post-rationalization \citep{wallat2024correctness}; our
claims are confined to correctness.

\section{Benchmark}
\label{sec:bench}

\paragraph{Corpus.}
50 public-domain FDA drug labels from DailyMed, fetched via
\texttt{evidence-eval fetch dailymed} (idempotent; manifest with
per-file SHA-256). Public domain, high-stakes, factually dense ---
labels carry crisp doses, concentrations, and contraindications.

\paragraph{Seeds.}
254 \texttt{valid} seed examples, each a self-contained world:
\texttt{corpus\_spans}, \texttt{prompt\_chunks} (what the model was
shown), \texttt{response\_sentences} (the fixture answer with
\texttt{cited\_spans}), and author-supplied
\texttt{support\_mutations}. Fact types are rotated (quantitative,
contraindication, indication, administration, multi-span). Authoring
provenance and limitations are documented in a datasheet
(\texttt{fda\_label\_bench\_PROVENANCE.md}).

\paragraph{Verification and the v1\,$\to$\,v2 audit.}
Structural validity is machine-checked by \texttt{evidence-eval lint}
(whole-corpus: 0 diagnostics). We report one methodological event
because it bears on validity. An initial benchmark (v1) showed a 49\%
false-refusal rate on valid examples; a hand audit against source PDFs
traced the cause to an authoring artifact (line-wrapped PDF spans
cited as fragments), not a property of the validator. We corrected the
guideline, re-authored the seeds (v2), and repeated the run; the
residual decomposed into ${\sim}15$ points fixed bug and ${\sim}29\%$
genuine NLI conservatism. Catching this by audit is itself part of the
evidence base (\texttt{benchmark-results.md}).

\paragraph{Failure injection (named operators, \citealp{gardner2020contrast}).}
From each valid seed, \texttt{evidence-eval inject} derives labeled
failures:
\begin{itemize}
  \item \emph{Existence:} rewrite a citation to a span ID absent from
    the corpus.
  \item \emph{In-context:} duplicate a span's text at a fresh ID
    outside the prompt, repoint the citation --- the duplicate text
    controls the relevance confound (a true matched pair: only
    in-context status changed).
  \item \emph{Support:} apply an author-supplied mutation labeled
    \texttt{unsupported} (off-topic substitution) or
    \texttt{contradicted} (number/negation/entity perturbation) on the
    same cited span.
\end{itemize}
254 seeds $\to$ 1{,}270 examples (254 valid + 1{,}016 injected).
Operators are named, deterministic, individually unit-tested.

\paragraph{Natural-failure supplement \texttt{[OPEN --- not run]}.}
Real LLM outputs (no injection), human-graded into the taxonomy,
defending against ``injected failures are too easy.'' Requires a
held-out corpus, a generation model, and human labeling ---
deliberately not automated.

\section{Experiments}
\label{sec:exp}

Four nested policies: \texttt{vanilla} $\subset$
\texttt{existence\_only} $\subset$ \texttt{two\_gate} $\subset$
\texttt{three\_gate}. Support gate evaluated under distilbert-MNLI and
DeBERTa-v3-MNLI.

\begin{itemize}
  \item \textbf{\S5.1 Model effect}
    (\texttt{nli-comparison.md}):
    distilbert\,$\to$\,DeBERTa moves three-gate agreement
    $1{,}106 \to 1{,}166/1{,}270$; valid-retention
    $167 \to 181/254$; structural classes $\Delta = 0$ (rules 1--2
    model-independent).
  \item \textbf{\S5.2 Per-rule isolation}
    (\texttt{rule-metrics.md}):
    existence F1 = \textbf{1.000}, in-context F1 = \textbf{1.000}
    (exact, model-independent); support F1 = 0.929 (recall 0.992,
    precision 0.873 --- the precision gap is the \S5.1 conservatism,
    not a missed failure).
  \item \textbf{\S5.3 Non-overlap} (the kernel): out-of-context
    accepted by existence-only $254/254$ (100\%); support failures
    accepted by two-gate $508/508$ (100\%). \textbf{Fully disjoint on
    the injected benchmark (100\% by the operational blindness
    measure)} --- each failure class is \emph{invisible} to the rules
    preceding its gate. This is a measured blindness property (OOC
    failures are invisible to the existence gate; support failures are
    invisible to the two-gate policy), an empirical property of the
    data --- \emph{not} an artifact of the nested-policy construction.
  \item \textbf{\S5.4 Composition}: F1 ladder
    $0.000 \to 0.400 \to 0.667 \to \mathbf{0.963}$; \textbf{additive
    lift $+56.3$ points} over the best non-composed baseline
    (conservative lower bound --- nested policies).
  \item \textbf{\S5.5 ALCE comparison \texttt{[OPEN]}},
    \textbf{\S5.6 natural-failure \texttt{[OPEN]}},
    \textbf{\S5.7 cost \texttt{[OPEN]}} --- not yet run; no fabricated
    numbers.
\end{itemize}

The decomposition is justified iff each rule isolates its class and
the rules don't overlap; both hold on the v2 injected FDA-label
benchmark with hard numbers. The structural core is exact and
model-independent; the one soft number is the separately-measured
support-gate conservatism.

\section{Limitations}
\label{sec:limits}

\begin{itemize}
  \item \textbf{Domain.} FDA labels only. No cross-domain
    generalization claim.
  \item \textbf{Size.} 254 seeds / 1{,}270 examples --- workshop-%
    appropriate, not main-conference scale.
  \item \textbf{Benchmark provenance.} Seeds are AI-assisted-authored
    from a fixed guideline and are therefore structurally more
    homogeneous than independent hand curation. We disclose this
    openly. The structural results (\S5.2 existence/in-context, all of
    \S5.3) are authoring-independent: existence and in-context are
    exact set operations whose outcomes do not depend on how seeds
    were written. Authoring homogeneity bears only on the support-gate
    numbers, and we flag it as a fair reviewer question. A sample was
    audited against source PDFs; an exhaustive human pass remains
    future work. (Datasheet:
    \texttt{fda\_label\_bench\_PROVENANCE.md}.)
  \item \textbf{NLI error propagation.} The support gate inherits the
    NLI model's errors; we measure NLI behavior as a separate column
    rather than bury it (the support precision figure \emph{is} that
    measurement).
  \item \textbf{Injection by construction.} \S5.6 is the intended
    defense and is not yet run.
  \item \textbf{No human-factors claim.} We do not measure user trust
    or downstream decision quality.
  \item We measure citation \emph{correctness}, not causal citation
    faithfulness; the in-context gate proves availability, not
    reliance (post-rationalization is possible). This is scoped
    explicitly in \S\ref{sec:rules}.
  \item Headline F1/lift numbers come from a class-balanced injected
    set where failures are prevalent; real RAG has lower failure
    prevalence where false-refusal cost dominates. We report the
    support gate's false-refusal rate explicitly rather than optimize
    it away.
\end{itemize}

\section{Conclusion and Future Work}
\label{sec:conc}

Citation correctness decomposes into three operationally independent
rules. On the v2 injected FDA-label benchmark --- a controlled
benchmark, not a real-world prevalence study --- the structural rules
are exact and model-independent, the rules catch fully disjoint error
classes (100\% by the operational blindness measure), and the
composition beats the best single rule by a wide margin. The semantic
support gate is the open frontier: it catches real failures but is
conservative, which we measure rather than hide.
Future work:
(1) larger, cross-domain, human-curated benchmark;
(2) clause-level support decomposition (FActScore-style) to attack the
support-precision ceiling;
(3) the natural-failure study (claim 4) and ALCE translation;
(4) production integration with user studies.
The system, benchmark, and one-command reproduction are released as
\texttt{evidence} (open source; version pinned to these experiments
--- see \texttt{PAPER\_EXPLAINER.md} and \S\ref{sec:conc} artifact
appendix).

\section*{Authoring Disclosure}

Portions of this manuscript and the benchmark seed authoring were
prepared with AI assistance from a fixed guideline; the lead author
designed the methodology, conducted the v1\,$\to$\,v2 audit, verified
all reported numbers against the released artifacts, and is
responsible for all claims. The structural results
(\S\ref{sec:exp}: existence, in-context, non-overlap) are exact set
operations whose outcomes do not depend on how seeds were written;
authoring homogeneity bears only on the support-gate numbers, as noted
in \S\ref{sec:limits}.

